\section{Standard tensor orbits with at least three components}
\label{sec:standard-orbit-cubic}

\paragraph{Accepted status and scope.}
D-TD52, question Q2fh, is a new derivation by
\texttt{/root/sixfolds}. Complete author, independent and root
whole-proof and packet readings pass. The new local ingredient is a cubic comparison for
the product of several standard character ratios, including the
one-fixed-component and two-fixed-component cases. The actual induced
covers and their genus argument are the accepted D-TD49 interface,
reproduced below. No source budget is assigned to the external partner.

Let \(C\longrightarrow\mathbf P^1\) be an actual connected finite
\(J\)-Galois cover, with a finite complex coefficient representation
\(V\) and genuine generating product-one inertia \(T_1,\ldots,T_r\).
Suppose
\[
 e=\dim V,\qquad
       \sum_{i=1}^r\operatorname{rank}(1-T_i)=2e+6.
 \tag{SOC.1}
\]
Suppose that the actual action has a stable whole tensor layer
\[
 V=\left(\bigotimes_{j=1}^t W_j\right)\otimes R,\qquad
 W_j\simeq\mathrm{Std}_m,\quad a=\dim W_j=m-1,\qquad
 m\ge2^{64},\quad 3\le t\le54m.
 \tag{SOC.2}
\]
Here \(R\) is the entire remaining space, of any positive dimension.
The represented core of this layer is \(A_m^t\), acting by its
irreducible standard modules. The original image permutes its component
matrix algebras transitively. Conjugation supplies an actual natural
wreath action
\[
             J\longrightarrow S_m\wr S_t
 \tag{SOC.3}
\]
whose image contains \(A_m^t\). These are hypotheses on the actual
extracted layer, not on an independently proposed lift to \(J\).
As in D-TD49, a coefficient-kernel element centralizes every represented
component and hence kills (SOC.3). The natural actions are therefore
actual quotient actions of the given cover.

Fix component identifications once. Compatible local normalizer
operators are standard \(S_m\)-operators up to scalars. Indeed their
conjugation automorphism lies in \(\operatorname{Aut}(A_m)=S_m\);
comparison with its standard representative differs on the perfect
core by a scalar character, which is trivial. Schur's lemma then makes
the remaining discrepancy scalar. A sign-twisted standard extension
has exactly the same absolute-character and projective-support bounds.

\begin{theorem}[D-TD52]
The actual source (SOC.1)--(SOC.3) does not exist. The theorem permits
an arbitrary entire partner \(R\), every scalar phase and every
coefficient field. It needs no twelve-mark premise, partner support
floor, maximality premise, or independent sphere source for a component.
\end{theorem}

\begin{proof}
\subsection*{1. Weighted traces and a cubic induced comparator}

For any actual finite-order coefficient operator \(g\), let \(\pi\)
be its component permutation, \(I=\operatorname{ind}(\pi)\), and \(c\)
the number of its fixed components. A weighted tensor cycle is conjugate
to cyclic rotation with one twist. Trace contraction around its cycles,
including the whole partner with normalized absolute trace at most one,
gives the accepted exact identity and inequality
\[
 \frac{|\operatorname{tr}(g)|}{e}
 \le a^{-I}\prod_{\mathcal C}
              \frac{|\operatorname{tr}(A_{\mathcal C})|}{a}
 \le a^{-I}\prod_{j:\,\pi j=j}\rho_j,\qquad
                 \rho_j=\frac{|f_j-1|}{a}.
 \tag{SOC.4}
\]
Here \(A_{\mathcal C}\) is the local cycle product, and \(f_j\)
is the number of fixed natural letters of the local permutation at
a fixed component. Each factor ratio lies in \([0,1]\).
Finite-order normalization is legitimate exactly as in ICL.6:
a power killing the component permutation is a tensor product of
the local cycle products and a partner power; a further power is
the identity, so each local factor has a scalar power. Rescaling makes
each factor finite order, and the remaining total scalar is a root
of unity. This applies separately to every original power of \(g\).
It does not replace its coefficient order by a natural permutation order.

Put \(\phi_j=(f_j/m)^3\). The elementary standard-character comparison is
\[
       \rho_j^3\le \phi_j+a^{-3}\mathbf1_{\{f_j=0\}}.
 \tag{SOC.5}
\]
For \(f_j\ge1\), this follows from
\((f_j-1)/(m-1)\le f_j/m\); for \(f_j=0\), \(\rho_j=1/a\).
Let \(\Omega\) be the disjoint union of the ordered cubes of the
\(m\) natural letters in all \(t\) components, allowing repetitions.
It is an actual \(J\)-set with
\[
 D=tm^3,\qquad F_\Omega(g)=\sum_{j:\,\pi j=j}f_j^3,\qquad
 \Phi(g)=\frac{F_\Omega(g)}D
                  =\frac1t\sum_{j:\,\pi j=j}\phi_j.
 \tag{SOC.6}
\]
We prove the uniform, all-power local estimate
\[
 \frac{|\operatorname{tr}(g)|}{e}\le\Phi(g)+\epsilon_t,\qquad
 \epsilon_t=
 \begin{cases}
       2a^{-3/2},&3\le t\le6,\\
       a^{-3},&7\le t\le54m.
 \end{cases}
 \tag{SOC.7}
\]

If \(c\ge3\), the arithmetic--geometric mean inequality gives
\(\prod\rho_j\le c^{-1}\sum\rho_j^c\le c^{-1}\sum\rho_j^3\).
Also \(a^{-I}\le c/t\). For \(\pi=1\) this is equality.
Otherwise put \(u=t-c\ge2\). Then
\(I\ge u/2\), and, since \(a\ge4\),
\(a^I\ge2^u\ge u+1\ge t/c\).
Consequently
\[
 \frac{|\operatorname{tr}(g)|}{e}
     \le\frac1t\sum_{j:\,\pi j=j}\rho_j^3
     \le\Phi(g)+\frac{c}{t a^3}\le\Phi(g)+a^{-3}.
 \tag{SOC.8}
\]
If every fixed component has \(f_j\ge1\), the error in (SOC.8)
is exactly zero.

If \(c=0\), (SOC.4) gives
\[
       \frac{|\operatorname{tr}(g)|}{e}\le a^{-I}
                         \le a^{-t/2}.
 \tag{SOC.9}
\]
This is at most \(a^{-3/2}\) when \(t\ge3\), and at most \(a^{-3}\)
when \(t\ge7\).

It remains to handle \(c=1,2\), where dropping the fixed-component
ratios would lose too much. Write \(A=a^{-I}\).
For every \(z\ge0\), the following exact optimization bounds hold:
\[
 Az-\frac{z^3}{t}
       \le\frac{2\sqrt t}{3\sqrt3}A^{3/2},\qquad
 Az^2-\frac{2z^3}{t}
       \le\frac{t^2}{27}A^3.
 \tag{SOC.10}
\]
For the first set \(u=\sqrt{At/3}\) and use
\((z-u)^2(z+2u)=z^3-3u^2z+2u^3\ge0\).
For the second set \(u=At/3\) and use
\((z-u)^2(2z+u)=2z^3-3uz^2+u^3\ge0\).
These prove (SOC.10) on the whole nonnegative half-line, so no
unproved location of a maximum in \([0,1]\) is needed.

For \(c=1\), substitute \(z=\rho_1\) and then (SOC.5).
For \(c=2\), put \(z=\sqrt{\rho_1\rho_2}\), so
\(\rho_1^3+\rho_2^3\ge2z^3\), and use the second inequality.
Thus
\[
 \frac{|\operatorname{tr}(g)|}{e}-\Phi(g)
 \le
 \begin{cases}
 \displaystyle\frac{2\sqrt t}{3\sqrt3}a^{-3I/2}
                            +\frac1{t a^3},&c=1,\\[4pt]
 \displaystyle\frac{t^2}{27}a^{-3I}
                            +\frac2{t a^3},&c=2.
 \end{cases}
 \tag{SOC.11}
\]

For \(3\le t\le6\), one has \(I\ge1\) in both cases.
In the first row, \(2\sqrt t/(3\sqrt3)\le2\sqrt2/3<1\),
so the error is less than \(2a^{-3/2}\).
The case \(c=2\) requires \(t\ge4\), because a permutation cannot
move exactly one point. Its error is at most
\[
          \left(\frac{36}{27}+\frac12\right)a^{-3}
                    =\frac{11}{6}a^{-3}<2a^{-3/2}.
 \tag{SOC.12}
\]

For \(t\ge7\), use \(a\ge7\) and \(t\le54m\le108a\).
If \(c=1\), then \(I\ge(t-1)/2\ge3\), and
\(2\sqrt t/(3\sqrt3)\le4\sqrt a\).
If \(c=2\), the integer \(I\) is at least three, and
\(t^2/27\le432a^2\).
The two rows of (SOC.11) are therefore at most
\[
 \left(\frac4a+\frac17\right)a^{-3}\le\frac57a^{-3},
 \qquad
 \left(\frac{432}{a^4}+\frac27\right)a^{-3}<a^{-3},
 \tag{SOC.13}
\]
respectively. For the second comparison it suffices that
\(432/7^4<3/7\).
Equations (SOC.8)--(SOC.13) prove (SOC.7) in every case.

\subsection*{2. Only ten original entries can bear an error}

Call an original inertia bad if it moves components, or if it fixes
all components and one local natural permutation has fewer than
\(m/2\) fixed letters. A nonbad inertia has only nonbad powers:
its component permutation is the identity, and each power fixes
every natural letter fixed originally. For all its powers,
\(c=t\ge3\) and every \(f_j\ge m/2\), so (SOC.8) has zero error.
This assertion includes original entries acting only through \(R\)
or a scalar.

Every bad inertia has normalized moving rank greater than \(1/5\).
For a component-moving inertia, the usual weighted-cycle word proof
gives support at least \((a-1)/(2a)>1/5\). More explicitly, on a
cycle of length \(\ell\ge2\), diagonalize its finite-order twist.
There are \(a\) constant words; every other cyclic word orbit has
length at least two and contributes at most one to a target eigenspace.
Its maximum multiplicity is at most
\((a^\ell+a)/2\), giving this support bound. Tensoring the other
cycles and the entire \(R\) preserves normalized projective support.

For a component-fixing bad inertia, let \(y\) be a local permutation
with \(f(y)<m/2\). On the natural permutation module every eigenspace
has dimension at most the number of cycles of \(y\). Passing to its
standard summand cannot increase that maximum. Hence
\[
 \frac{\nu(\mathrm{Std}_m(y))}{a}
 \ge\frac{\operatorname{ind}(y)-1}{m-1}
 >\frac{m/4-1}{m-1}>\frac15.
 \tag{SOC.14}
\]
The last inequality holds for \(m>16\). Again convolution transfers
this lower bound to the actual whole operator, independently of scalars.
Since \(e\ge a^t\ge a^3\) and \(a\ge7\),
\(2+6/e<11/5\). Eleven bad original entries would contradict (SOC.1).
There are therefore at most ten.

For an original coefficient order \(L_i\), cyclic Fourier inversion
at every target eigenvalue and (SOC.7) yield
\[
 \frac{\operatorname{rank}(1-T_i)}e
 \ge\frac{\nu(T_i)}e
 \ge\frac{\operatorname{ind}_\Omega(T_i)}D
               -\epsilon_t\mathbf1_{\{i\ {\rm bad}\}},\qquad
 \frac1e\sum_i\operatorname{rank}(1-T_i)
 \ge\frac1D\sum_i\operatorname{ind}_\Omega(T_i)-10\epsilon_t.
 \tag{SOC.15}
\]
The average is over all \(L_i\) original powers. The order of the
actual permutation on \(\Omega\) divides \(L_i\), so its average
fixed-point count is exactly its number of cycles. Identity natural
images have zero index and zero error, even when the original operator
is genuine. No projected mark count has been inserted.

\subsection*{3. The actual induced cubes in every base genus}

We include the geometric argument to specify its full scope.
The actual action on \(\Omega\) has five orbits: the diagonal,
three equality patterns with exactly two distinct letters, and the
ordered triples of distinct letters. The group \(A_m^t\) is
three-transitive within each component, and the component action is
transitive. Let \(Y_1,Y_2,Y_3\) be the resulting connected quotient
curves; the three two-distinct equality-pattern covers are isomorphic.
Their exact Riemann--Hurwitz identities give
\[
 \sum_i\operatorname{ind}_\Omega(T_i)
 =2D-10+2\bigl(g(Y_1)+3g(Y_2)+g(Y_3)\bigr).
 \tag{SOC.16}
\]

Let \(H\le J\) stabilize one component, and put \(X=C/H\),
\(h=g(X)\). The actual component homomorphism \(H\to S_m\)
has image \(P=A_m\) or \(S_m\), since it contains \(A_m\).
Its kernel \(K\) gives a connected \(P\)-Galois cover
\(C/K\longrightarrow X\). Its ordered-distinct-triple quotient
is \(Y_3\longrightarrow X\), of degree \((m)_3\).
In particular \(X\) is not assumed rational.

If \(h=0\), accepted HD-GC06 applies, since \(m>10^6\).
The unordered-three-set quotient \(Z\) has genus \(g(Z)>m^2/288\).
The forget-order map \(Y_3\to Z\) has degree six, including for
\(A_m\): parity on three selected letters can be matched on the
remaining letters. Thus \(g(Y_3)\ge6g(Z)-5>m^2/48-5\).
If \(h\ge2\), Riemann--Hurwitz gives \(g(Y_3)\ge1+(m)_3\),
which is stronger.

If \(h=1\), the component cover must ramify, since a finite connected
unramified cover of a complex elliptic curve has abelian monodromy.
A nonidentity branch permutation \(y\) fixes at most \(m-2\) letters.
The number of moved ordered distinct triples is at least
\((m)_3-(m-2)(m-3)(m-4)=6(m-2)^2\).
Its index is at least half that number. Hence
\[
       g(Y_3)\ge1+\frac32(m-2)^2>\frac{m^2}{48}-5.
 \tag{SOC.17}
\]
This proves the same quadratic lower bound in every component-base
genus. Together with nonnegativity of the other genera, (SOC.16) gives
\[
 \frac1D\sum_i\operatorname{ind}_\Omega(T_i)
       >2+\frac1{24tm}-\frac{20}{tm^3}.
 \tag{SOC.18}
\]
This is exactly the accepted ICL.14--17 geometry, with all original
cover maps and all their possible base genera retained.

\subsection*{4. Explicit numerical contradiction}

Since \(a=m-1\ge m/2\), one has
\[
                   e\ge a^3\ge m^3/8,\qquad
                   6/e\le48/m^3.
 \tag{SOC.19}
\]
For \(3\le t\le6\), equations (SOC.7), (SOC.15) and (SOC.18) imply
\[
 \frac1e\sum_i\operatorname{rank}(1-T_i)
 >2+\frac1{144m}-\frac7{m^3}-\frac{60}{m^{3/2}}
 >2+\frac1{288m}>2+\frac6e.
 \tag{SOC.20}
\]
Here \(10\epsilon_t\le20a^{-3/2}<60m^{-3/2}\).
At \(m\ge2^{64}\),
\(60/\sqrt m+7/m^2\le60/2^{32}+7/2^{128}<1/288\).
For instance each summand is less than \(1/576\).
Also \(m^2>48\cdot288=13824\).
These prove the last two strict inequalities in (SOC.20).

For \(7\le t\le54m\), the same equations yield
\[
 \frac1e\sum_i\operatorname{rank}(1-T_i)
 >2+\frac1{1296m^2}-\frac{83}{m^3}
 >2+\frac1{2592m^2}>2+\frac6e.
 \tag{SOC.21}
\]
Indeed \(10\epsilon_t=10a^{-3}\le80/m^3\) and
\(20/t<3\). The last comparisons follow from
\(m>83\cdot2592=215136\) and
\(m>48\cdot2592=124416\), both immediate at \(2^{64}\).
Both cases contradict (SOC.1), proving D-TD52.
\end{proof}

\paragraph{Exact near-maximum application and remaining boundary.}
Fix one complete D-TD12 recursion in the original faithful primitive
rank-nine source, let its largest alternating natural degree be
\(n\ge N_*\), and let an actual ordinary standard layer have
\(m\ge\lceil(n+120)/108\rceil\). Here \(N_*\) is the unchanged
HD-ML01/HD-TW03 integer. The accepted D-TD16 whole-orbit theorem gives
\[
 t\le\left\lfloor\frac{n-1}{2}\right\rfloor<54m,\qquad
          m>241^{482}>2^{64}.
 \tag{SOC.22}
\]
Its full original recursion and twelve genuine marks are supplied in
this application, not added to the standalone theorem.
Thus D-TD52 excludes every such standard layer with \(t\ge3\).
Together with D-TD49 and D-TD51 in their exact scopes, this leaves
standard layers with one or two components. This statement does not
assert that a two-component layer is stable componentwise, and it
does not apply a separate rank-six or rank-nine budget to any factor.
It constructs no source, polarized factor, period map, or detector.

\paragraph{Known inputs and source boundary.}
The standard-character identity is the elementary decomposition of the
natural permutation module into its trivial and sum-zero summands.
The weighted-cycle trace, finite-order normalization and actual induced
cover geometry are the complete accepted ICL.6--17 arguments, reproduced
here. The only genus theorem invoked is accepted HD-GC06 with its actual
sphere-cover hypothesis and numerical threshold \(10^6\).
The \(c=1,2\) polynomial optimization and the uniform error estimates
above are new and proved in full. No new external classification result,
unrefereed theorem, unevaluated asymptotic constant, or assumption on the
entire partner has been introduced.
